\cotefile{36}
\scanpage{3}
\mspage{1}

% head
\uline{Remarques.\qquad SGA 7 IX}

% upper left
0.

% i
i).\quad Dire que $l$ est un nombre premier à $p$\quad (pour que $\begin{array}[t]{c}l\,6\\l\,7\\\mathrm{ii)}\,l\,3\end{array}$ soit ok p.e.)

% i, second line
Pourquoi $S$ hensélien ?

% 0.1
0.1, $l\,9$ : pourquoi seulement à isogénie près ? (si $l\ne p$;\deleted{\unsure{)}} sinon ce ne veut rien dire)

% 0.1, next line
$l\,14$ : il semble clair \qquad est un peu optimiste.

% ii
ii).\quad $l\,14$\qquad français :\qquad et que \qquad ( ; il\qquad ? )

% iii
iii).\quad \tikz[baseline=(c.base)]{\node[draw,ellipse,inner sep=1pt] (c) {$l\,6$};}\qquad une référence ou définition pour réduction semi-stable d'une courbe (seule la terminologie de ``réduction stable'' se trouve dans $[3]$)

% iv
iv)\quad $l\,12$\qquad orthographe ?\qquad système

% iv, next line
\tikz[baseline=(c.base)]{\node[draw,ellipse,inner sep=1pt] (c) {$l\,15$};}\qquad vider la Rmq sur les $H^i$ (qui n'est ni \tikz[baseline=(c.base)]{\node[draw,ellipse,inner sep=1pt] (c) {liée à la ?};}monodromie, ni fondé sur des théorèmes)

% division
\hrule\smallskip\hrule\smallskip\hrule

% lower section
1.

% p 0, title
\uline{p 0}. titre :\qquad $l'$\quad (minuscule)

% 1.0
(1.0) $l\,5$\qquad donc\qquad ?

% 1.0, next line
\tikz[baseline=(c.base)]{\node[draw,ellipse,inner sep=1pt] (c) {$l\,7$};}\qquad à quel concept se rapporte ``lequel''.

% division
\hrule

% p 1
\uline{p 1}\quad $l\,12$\qquad français : choisie \qquad ; ``qui est'' tend à faire croire qu'on impose une condition supplémentaire

% p 1, next line
$l\,16$\qquad $=B$ dans le symbole est malheureux ;\quad $=A'_K$ de toute façon.

% division
\hrule

% p 2
p 2\quad $l\,2$\qquad pourquoi \tikz[baseline=(c.base)]{\node[draw,ellipse,inner sep=1pt] (c) {G};}roupe\qquad -- épimorphisme n'est pas la définition d'isogénie --

% p 2, l 15
$l\,15$\qquad -- supprimer séparé type fini au bénéfice du seul quasi projectif.

% last line
-- que $A\mathbin{\mathop{\times}\limits_S}\operatorname{Spec}(K)\overset{\sim}{\longrightarrow}A_K$ est une trivialité
% archive sheet number: 2

\scanpage{4}
\mspage{2}

% opening
$l-8$\qquad pourquoi \uline{donc}\qquad ?

% division
\hrule

% p 3
p 3\quad $l-8,-9$\qquad $B$ ?

% division
\hrule

% p 4, 1.3.1
p 4.\quad 1.3.1.\qquad As tu vérifi\deleted{\unsure{e}}é si l'accouplement d'A. M. coïncide avec le tien ?

% p 4, l 4
$l\,4$.\qquad \uline{formellement} non ramifié ?

% p 4, l 7
$l\,7$.\qquad \uline{séparablement} clos ?

% p 4, l 14–20
$l\,14-20$.\quad Est-ce utile à dire, en l'absence de référence à ce que ce peut vouloir dire ?

% p 4, l −5
$l-5$\qquad \unsure{Qu'il|qu'il} soit \qquad et \qquad signifierait \qquad plutôt ?

% division
\hrule

% pg 5
pg 5\quad $l-2$\qquad français :\quad les classes $\unsure{w|\omega}$\qquad \uline{ou}\qquad relatifs.

% division
\hrule

% pg 6
pg 6\quad $l\,4$\qquad français. dire\qquad un Module inversible ample\qquad ?

% pg 6, below
\deleted{\illegible{4}}

% division
\hrule\smallskip\hrule\smallskip\hrule

% section
\tikz[baseline=(c.base)]{\node[draw,double,ellipse,inner sep=1pt] (c) {2};}

% pg 7
pg 7\qquad français \uline{$l-9$}\quad \tikz[baseline=(c.base)]{\node[draw,ellipse,inner sep=1pt] (c) {$\unsure{c|\subset}$};}$L_0$\quad donc \tikz[baseline=(c.base)]{\node[draw,ellipse,inner sep=1pt] (c) {c'};} est

% pg 7, second line
\deleted{$l- \unsure{8|7|1}$ \unsure{si} \unsure{que}}\quad $l-7$\quad \tikz[baseline=(c.base)]{\node[draw,ellipse,inner sep=1pt] (c) {que};}\quad (le $2^{\underline{\mathrm o}}$)

% division
\hrule

% pg 8
pg 8\qquad français\quad $l\,4$\qquad par exemple \inserted{\tikz[baseline=(c.base)]{\node[draw,ellipse,inner sep=1pt] (c) {pour};} ?}\unsure{k|si} $\ldots$

% pg 8, next line
$l-7$\quad \uline{le} terme

% division
\hrule

% pg 14
pg 14.\qquad L'argument est un peu immoral ; si $X$ et $Y$ sont extension d'une VA par un tore et si $\varphi:X\longrightarrow Y$ est donné, alors $\varphi$ est strictement compatible à la filtration par le poids et $\varphi$ est une isogénie si et seulement si $T_\ell(\varphi)$ en est une.

% division
\hrule

% pg 20
\tikz[baseline=(c.base)]{\node[draw,ellipse,inner sep=1pt] (c) {pg \deleted{\unsure{15|18}} 20.};}\quad Rmq 2.6 a)\qquad $S$ est-il un trait dans cette remarque ; sait-on assez sur le passage à la limite et Néron ?
% archive sheet number: 3

\scanpage{5}
\mspage{3}

% section 3, pg 23
\tikz[baseline=(c.base)]{\node[draw,ellipse,inner sep=1pt] (c) {3};}\quad \uline{pg 23}\qquad On peut dire plus simplement que $\,{}_nK$ est étale séparé car $\,{}_nK=\operatorname{Ker}({}_nG\longrightarrow{}_nH)$. En a), ne dirait-tu pas que $\,{}_nK_{\bar s}$ est isomorphe à un sous-groupe de $\,{}_nK_{\bar t}$.

% pg 24
\uline{pg 24}\quad d) $l\,15$ :\quad \tikz[baseline=(c.base)]{\node[draw,ellipse,inner sep=1pt] (c) {$\dim$};}$\operatorname{Ker}(u_{\unsure{s|1}})=0$ ?

% pg 28
\tikz[baseline=(c.base)]{\node[draw,ellipse,inner sep=1pt] (c) {\uline{pg 28}};}\quad Les notions introduites dans 3.5 ont elles été définies pour $S$ non hensélien ? (il faudrait au moins choisir un groupe d'inertie).

% pg 28, below
[item pour 3.7]

% pg 29
\uline{pg 29}\quad $l-2$\quad la terminologie ``clôture normale'' est-elle encore canonique ?

% division
\hrule\smallskip\hrule\smallskip\hrule

% section 4, p 32
\tikz[baseline=(c.base)]{\node[draw,ellipse,inner sep=1pt] (c) {4};}\quad \uline{p 32}\quad $l\,8$\qquad ne dépend \tikz[baseline=(c.base)]{\node[draw,ellipse,inner sep=1pt] (c) {pas};}

% pg 33 ss
\tikz[baseline=(c.base)]{\node[draw,ellipse,inner sep=1pt] (c) {? \quad\uline{pg 33 ss}};}\quad (4.2)\qquad Ne donnerais tu pas un énoncé style ``voisinage tubulaire'' compréhensible. la vérité est

% pg 33 ss, continuation
\deleted{\illegible{1}} on a un morphisme de topos $\operatorname{Spec}(K)_{\mathrm{et}}\xrightarrow{\varphi}\operatorname{Spec}(k)_{\mathrm{et}}$ (``\unsure{ent.|ext.} du corps résiduel'') et $A_0^0$ est un schéma relatif sur $(\operatorname{Spec}(K))_{\mathrm{et}}$ annelé par $\varphi^*\unsure{O|\mathcal O}$ [le $\varphi^*$ des faisceaux d'ensembles] -- Même chose d'ailleurs pour $\underset{\text{ramifier.}}{\varinjlim}\, A_0$

% pg 36
\tikz[baseline=(c.base)]{\node[draw,ellipse,inner sep=1pt] (c) {si ?\quad\uline{pg 36}};}\quad $l-6$.\qquad \unsure{*|``} ce n'est pas la trace qui est contragrédiente ; l'énoncé primitif était correct --

% pg 40
\uline{pg 40}\qquad (4.5.4)\quad Qu'est-il écrit dans la marge ?

% pg 40, l −1
$l-1$\qquad Quel est loc.cit. ?

% pg 40, l 3
$l\,3$\qquad $l$-groupe signifie-t-il tué par $l$, ou par un $l^k$ ?

% last line
le trait est-il strictement local ? ce serait plus agréable.
% archive sheet number: 4

\scanpage{6}
\mspage{4}

% pg 41
\tikz[baseline=(c.base)]{\node[draw,ellipse,inner sep=1pt] (c) {\uline{pg 41}};}\quad $l-2$\qquad ``fonction de Dirac'' est malheureux (sur un groupe, on ne sait trop si c'est \quad ou $\left\{\begin{array}{l}f(e)=1,\;0\ \text{ailleurs}\\f(e)=\#G,\;0\ \text{ailleurs}\end{array}\right.$ )

% left of pg 41
% slanted
\margin{guillemets suffisent ?}

% pg 42
\uline{pg 42}\quad début dém. 4.7 : le \unsure{cas|aps} n'est pas clair pour $n=2$, et je n'ai pas confiance en dehors du cas strictement local.

% division
\hrule\smallskip\hrule\smallskip\hrule

% section 5
\tikz[baseline=(c.base)]{\node[draw,ellipse,inner sep=1pt] (c) {5};}

% pg 45
\uline{pg 4\deleted{\unsure{7}}5} $l\,3$ ss\quad français : beaucoup de ``donc''.

% pg 46
\uline{pg 46}\quad Quel est le sens de (5.2.3,1) et de \deleted{\unsure{le|la}} ``par définition'' qui suit -- S'agit-il de pro-objets, ou de schémas.

% left of pg 46
\margin{\tikz[baseline=(c.base)]{\node[draw,ellipse,inner sep=1pt] (c) {je le dis !};}}

% pg 48
\uline{pg 48}\quad $l-4$\quad Pourquoi as-tu barré connexe ?

% pg 53
\deleted{\uline{pg 53}\quad 5.7 est une redite pour $l\ne p$, et un peu idiot pour $l=p$.}\quad (utilisé en 5.15)

% pg 56 bis
\tikz[baseline=(c.base)]{\node[draw,ellipse,inner sep=1pt] (c) {?\quad\uline{pg 56 bis}\quad 5.15.};}\quad tu n'a plus besoin ici de $\operatorname{car}(K)\ne l$

% 54 quater
54 quater -- 5.17.\qquad seraient plutôt que seront.

% 55 ter
\uline{55 ter}\quad $l-10$. Il serait honnête de dire que 6.6 utilise $\operatorname{car}(K)=0$, via Tate

% left of 55 ter
% slanted
\margin{\tikz[baseline=(c.base)]{\node[draw,ellipse,inner sep=1pt] (c) {\shortstack{c'est dit\\dans les \unsure{hyp.|hyp}\\de 5.13}};}}

% division
\hrule\smallskip\hrule\smallskip\hrule

% section 6, pg 57/58
\tikz[baseline=(c.base)]{\node[draw,ellipse,inner sep=1pt] (c) {6};}\quad \uline{pg 57/58} : tu supposes tantôt $A^0$ $l$-divisible (58, $l-11$), tantôt non (57, $l\,3$)

% pg 58
\uline{pg 58}\quad $l-5$\quad $({}_{p^\infty}X)^{\text{\tikz[baseline=(c.base)]{\node[draw,ellipse,inner sep=1pt] (c) {0};}}}$\quad \tikz[baseline=(c.base)]{\node[draw,ellipse,inner sep=1pt] (c) {?};}\qquad itou plus loin --

% pg 59
\uline{pg 59}\qquad Je ne comprends rien -- Que fais-tu de la partie discrète ? --

% left of pg 59
% slanted
\margin{\uline{CANULAR ?!}}

% pg 59, continuation
Je ne crois pas que la plus grande \deleted{\unsure{s|ss|de}} sous VA ordinaire puisse se relever (\~m en cas de bonne réduction).

% last lines
\uline{i e}\quad 63 est faux\quad (déformer de façon générale $E\times F$, $E$ et $F$ elliptiques, $E$ ordinaire, $F$ non ordinaire ; général $\Rightarrow$ générique simple.
% archive sheet number: 5

\scanpage{7}
\mspage{5}

% pg 60
\uline{pg 60}\qquad Je n'ai pas confiance en ``il est immédiat'' de 6.4.

% left of pg 60
% slanted
\margin{\tikz[baseline=(c.base)]{\node[draw,ellipse,inner sep=1pt] (c) {\shortstack{c'est pourtant\\trivial.}};}}

% pg 61, 61 bis
\uline{pg 61, 61 bis} : Je n'ai pas lu.

% next paragraph
Il me semble qu'il doit être possible de définir, à isogénie près, des ``parties finies'' et ``parties toriques'', ayant le bon rang --

% left of preceding paragraph
% slanted
\margin{\tikz[baseline=(c.base)]{\node[draw,ellipse,inner sep=1pt] (c) {\shortstack{non, mon\\argument\\marche}};}}

% pg 64
\uline{pg 6\deleted{\unsure{3|5}}4/}\deleted{\unsure{6.6|66|68}} : \deleted{\unsure{$l$ 10|$l$ 11|Serre}} (réf. à Serre pour l'indép. de $l$ ligne 11 ?)

% pg 64, second line
Un peu malhonnête d'être si bref, car 6.6 est utilisé en 5.13.

% left of pg 64
% slanted
\margin{\tikz[baseline=(c.base)]{\node[draw,ellipse,inner sep=1pt] (c) {que faire ?};}}

% division
\hrule\smallskip\hrule\smallskip\hrule

% section 7, p 65
\tikz[baseline=(c.base)]{\node[draw,ellipse,inner sep=1pt] (c) {7};}\quad \uline{p 65}.\qquad $A_0$ n'est pas défini\qquad \margin{si}

% p. 67
\margin{\tikz[baseline=(c.base)]{\node[draw,ellipse,inner sep=1pt] (c) {?};}}\uline{p. 67}\qquad (7.1.4)\qquad français --

% division
\hrule\smallskip\hrule\smallskip\hrule

% section 8
\tikz[baseline=(c.base)]{\node[draw,ellipse,inner sep=1pt] (c) {8};}

% division
\hrule\smallskip\hrule\smallskip\hrule

% section 9, p 80
\tikz[baseline=(c.base)]{\node[draw,ellipse,inner sep=1pt] (c) {9};}\quad \uline{p 80}\qquad 9.2.3\qquad $u$ plutôt que $u(g)$

% p 82
\uline{p 82}\qquad (9.3.3)\qquad $+$ compatibilité avec $R$ ? --

% p 87
\uline{p 87}\qquad 9.4.\qquad $S$ \uuline{strictement} hensélien ? (on le suppose en p 88, $l-9$)

% next indented line
pas d'hyp. d'excellence ?

% p 91
\uline{p 91} :\quad 9.5.8.\qquad Ne mettrait-tu pas une génuflexion ``au signe près''.

% pg 92 bis
\uline{pg 92 bis}\qquad $S$ est-il strictement hensélien ?

% division
\hrule\smallskip\hrule\smallskip\hrule

% section 10, pg 98
\tikz[baseline=(c.base)]{\node[draw,ellipse,inner sep=1pt] (c) {10};}\quad \uline{pg 98}\qquad 10.5.3.\quad Une extension de $K$ (qui ici est infini) est-elle requise ?? .

% last line
10.5.4.\qquad Dire \uline{où} on fera ces vérifications --
% archive sheet number: 6

\scanpage{8}
\mspage{6}

% section 11, pg 101
\tikz[baseline=(c.base)]{\node[draw,ellipse,inner sep=1pt] (c) {11};}\qquad pg 101\qquad $l-6$.\qquad Dire où on trouve que Néron commute à $S\mapsto S^{\wedge}$.

% pg 103
pg 103\qquad Le renvoi au courage du lecteur est désagréable --

% left of pg 103
% slanted
\margin{\tikz[baseline=(c.base)]{\node[draw,ellipse,inner sep=1pt] (c) {que faire ?};}}

% last paragraph
Si j'ai bien compris, tu affirme que prolonger une \unsure{bientension|biextension} à $A\times A'$ est ``la \~m chose'' que de la prolonger aux $\,{}_{\ell^\infty}A\times{}_{\ell^\infty}A'$.
% archive sheet number: 7

\scanpage{28}

% upper left, two pen traces
\figure{\begin{tikzpicture}[x=0.26cm,y=0.26cm]
\draw[line width=.28pt] (0.3,2.4) .. controls (0.25,1.8) and (0.05,0.8) .. (0.05,0.55) .. controls (0.1,0.28) and (0.35,-0.12) .. (0.48,0);
\draw[line width=.22pt] (1.25,3.1) -- (1.1,1.9);
\end{tikzpicture}}

% upper right
\diagram{\begin{tikzcd}[ampersand replacement=\&,column sep=1.15em,row sep=.45em]
\& M\arrow[dr] \& \text{Situation} \& \& \& \\
\unsure{\cdot|0}\arrow[r,dash] \& T\arrow[r]\arrow[d,equal] \& G\arrow[r,dash] \& B \& 0\arrow[l] \& \\
\& D(M') \& \& \times \& \& \unsure{1|,}\\
\cdot\arrow[r,dash] \& T'\arrow[r]\arrow[d,equal] \& G'\arrow[r,dash] \& B'\arrow[r] \& 0 \&\\
\& D(M) \& \& \& \&\\
\& M'\arrow[uur] \& \deleted{\unsure{M|R|j}} \& \& \&
\end{tikzcd}}
% archive sheet number: 27

\scanpage{29}

% head, left
$\Gamma=\mathbb G$

% head, right
$M\times B'\qquad \operatorname{Ext}(B',DM)$

% head, second line
$\deleted{\operatorname{Hom}(B',\operatorname{Ext}(\unsure{M|B},\unsure{G|A|\mathbb G_m}))}=0$

% upper left and centre
\diagram{\begin{tikzcd}[ampersand replacement=\&,column sep=1.6em,row sep=1.7em]
B \& B' \& \text{une biext. }\inserted{\mathcal E}\text{ vers }\mathbb G_m\\
M\arrow[u,"i"] \& M'\arrow[u,"i'"'] \&
\end{tikzcd}}

% upper pair of maps
\[
\begin{aligned}
 B'&\xrightarrow{\alpha'}\underline{\operatorname{Ext}}(B,G_m)\mathbin{-}\underline{\operatorname{Ext}}(M,G_m)\quad\text{nul}\\
 B&\mathrel{\overset{\alpha}{\text{---}}}\underline{\operatorname{Ext}}(B',G_m)\mathbin{-}\underline{\operatorname{Ext}}(M',G_m)\quad\text{nul}
\end{aligned}
\]

% right of upper maps
\margin{$\begin{array}{l}\text{p. ex }\operatorname{Hom}(B',\underline{\operatorname{Ext}}(M,\overset{\substack{0\\\Vert}}{G_m}\\ \text{p.ex }\operatorname{Hom}(B,\underline{\operatorname{Ext}}(M',\underset{\substack{\Vert\\0}}{G_m}\end{array}$}

% upper middle
d'où extensions

% two extension rows
\diagram{\begin{tikzcd}[ampersand replacement=\&,column sep=1em,row sep=.35em]
\cdot\arrow[r,dash] \& DM\arrow[d,equal]\arrow[r,"j'"] \& G'\arrow[r,"p'",dash] \& B'\arrow[r] \& 0\\
\& T' \& \& \&\\
0\arrow[r] \& DM'\arrow[d,equal]\arrow[r,"j",dash] \& G\arrow[r,"p",dash] \& B\arrow[r] \&0\\
\&T\&\&\&
\end{tikzcd}}

% right of extensions, upper
\margin{$\mathcal E_{M,G'}$ est canoniqt splitté par $\rho$}

% right of extensions, lower
\margin{$\mathcal E_{G,M'}$ est can. splitté $\rho'$}

% example
\uline{Exemple}\qquad Si \deleted{\unsure{$B$ est}} $\underline{\alpha'}$ \deleted{\unsure{éta}}\inserted{est} un isom, alors $i'$ \unsure{revient à donner|revient : la donnée|revient au même de donner} un $i'\in\operatorname{Hom}(M',\underline{\operatorname{Ext}}(B,G_m))$

% example, arrows and conditions
\diagram{\begin{tikzcd}[ampersand replacement=\&,column sep=.5em,row sep=.45em]
{}\& \arrow[l,shift left=.4ex]\arrow[r,shift left=.4ex] \& \operatorname{Biext}(B,M';G_m) \& \arrow[l,shift left=.4ex]\arrow[r,shift left=.4ex] \& \deleted{\unsure{\operatorname{repr}}}\inserted{\operatorname{Ext}(B,D(M'))}\\
\&\text{si }\deleted{\unsure{\operatorname{Hom}}}\operatorname{Ext}(M',\underline{\operatorname{Hom}}(B,G_m))=0 \&\& \text{si }\operatorname{Hom}(B,\underline{\operatorname{Ext}}(M',G_m))=0 \&\\
\&\text{p.ex }\underline{\operatorname{Hom}}(B,G_m)=0\&\&\text{p.ex. }\underline{\operatorname{Ext}}(M',G_m)=0\&
\end{tikzcd}}

% middle, continuation
aussi \unsure{bien|pour|donc}. Donc la donnée de $i'$ revient à celle de l'ext. $G$. Et réciproquement

% lower half
Considérons l'image inverse \inserted{$\mathcal E_{\unsure{i,i'|i',i'}}$} de $\mathcal E$ sur $(M,M')$. Alors des splittings \inserted{$\rho_i$} de cette biextension correspondent aux hom $j$ de $M$ dans $G$ relevant $i$, ou aux \unsure{mor}$\,j'$ de $M'$ dans $G'$ relev. de $i'$.

% interlinear diagram
\diagram{\begin{tikzcd}[ampersand replacement=\&,column sep=1.4em,row sep=.2em]
 M\arrow[r,dash]\& G\arrow[r,dash]\& B\\
 M'\arrow[r,dash]\&G'\arrow[r,dash]\&B'
\end{tikzcd}}

% above cancellation
\deleted{\unsure{sur le}}

% lower text and rewritten block
Des relèvements \deleted{\unsure{Si}} $j$ et $j'$ \deleted{\unsure{se}} correspondent ainsi, \deleted{alors \unsure{on trouve un} splitting canonique \deleted{\unsure{de}} $\sigma$ de $\mathcal E_{G,M'}$ et un splitting canonique $\sigma'$ de $\mathcal E_{M,G'}$, ces deux splittings coïncident sur $M,M'$ \unsure{avec le} splitting \unsure{donné|associé} $h$ de $\mathcal E_{M,M'}$.}

% last three lines
ssi les deux splittings de $\mathcal E_{M,M'}$, déduits d'une part du splitting $\rho$ de $\mathcal E_{M,G'}$ et de $j':M'\mathbin{-}G'$, d'autre part du splitting $\rho'$ de $\mathcal E_{G,M'}$ et de $j:M\mathbin{-}G$, coïncident.
% archive sheet number: 28

\scanpage{30}
\begingroup\csname tagsleft@true\endcsname

% top, first member
or le premier membre (grâce à (2.3.2)) est $\phi^{n'}(x,my)=m\,\phi^{n'}(x,y)$. La compatibilité \deleted{(2.3.1)} (2.\deleted{2}3.2) \deleted{est}\inserted{implique} à son tour \deleted{conséquence de} la compatibilité symétrique\inserted{ de (2.3.1)}, exprimée par la commutativité de

% upper square, label at left
\diagram{\begin{tikzcd}[ampersand replacement=\&,column sep=small]
\& \deleted{\unsure{n'|m|\mathrm{id}}}\;{}_nP\arrow[d,"\handwritten{\mathrm{id}}"] \& \mathsf{x} \& {}_{n'}Q\arrow[d,"\handwritten{m}"]\arrow[rr,"{\phi_{\handwritten{\xi}}^{n'}\handwritten{|}\,{}_nP\;\mathsf{x}\;{}_{\handwritten{n'}}Q}"] \& \& {}_nG\arrow[d,"\handwritten{\mathrm{id}}"]\\
\text{(2.3.\deleted{\unsure{3}}4)} \& {}_nP \& \mathsf{x} \& {}_n\handwritten{Q}\arrow[rr,"{\phi_{\handwritten{\xi}}^n}"] \& \& {}_nG\,, 
\end{tikzcd}}

% below square
qui s'écrit en formules

% upper formula
\[\phi^n(x,my)=\phi^{n'}(x,y)\qquad (x\in{}_nP(S),\ y\in{}_nQ(S)).\tag*{(2.3.\deleted{\unsure{4}}5)}\]

% middle, crossed opening and reduction
\deleted{Montrons comment (2.3.\unsure{2|4}) implique (2.3.\unsure{4|1|3}) : avec les notations de cette dernière formule,} Comme $\,{}_{mn'}P\handwritten{\xrightarrow{\mathrm m}}{}_nP$ est un épimorphisme, on est ramené à prouver (2.3.2) dans le cas où $x=mx'$, où $x'\in{}_{n'}P(S)$, \deleted{donc}\inserted{or} on a $\deleted{\phi^n(mx,y)=\phi^n(m^{\inserted{\handwritten{2}}}x',y)=\phi^n}\quad\phi_{\handwritten{\xi}}^n(mx',my)\overset{\inserted{\text{(2.3.3)}}}{=}m\,\phi_{\handwritten{\xi}}^{n\handwritten{'}}(x',y)=\phi_{\handwritten{\xi}}^{n'}(mx',y)$ (par bilinéarité de $\phi^{n'}$) $=\phi_{\handwritten{\xi}}^{n'}(x,y)$ puisque $x=mx'$ \deleted{\unsure{q|g}}dans $\,{}_{n'}P(S)$, ce qui établit (2.3.5).

% lower middle, functoriality
Prouvons donc la compatibilité (2.3.1). La fonctorialité\inserted{ en $Q$} de l'homomorphisme (2.1.2) nous montre qu'on peut supposer $nQ=0$. Alors (2.1.2) est le composé

% lower middle, composite
\[\mathrm{Biext}^1(P,Q;G)\handwritten{\xrightarrow{\sim}}\mathrm{Hom}(\underline{\mathrm{Tor}}_1(P,Q),G)\handwritten{\xleftarrow{\sim}}\mathrm{Hom}({}_nP\otimes Q,G)\,,\]

% lower middle, cancelled line
\deleted{en $Q$ nous montre que (2.1.2) est le composé}

% lower middle, cancelled display
\[\deleted{\mathrm{Biext}^1(P,Q;G)\qquad\mathrm{Hom}(\mathrm{Tor}_1(P,Q),G)\qquad\mathrm{Hom}(\mathrm{Tor}_1(P,{}_nQ),G)}\]

% lower, two arrows
où la première flèche est celle de (1.1.3) et ne dépend pas du choix de $n$, et où la deuxième flèche est transposée de l'isomorphisme $\underline{\mathrm{Tor}}_1(P,Q)\handwritten{\xrightarrow{\overset{\unsure{\partial|d}}{\sim}}}{}_nP\otimes Q$ défini dans (1.3.4). Donc la compatibilité voulue équivaut essentiellement à \deleted{celle} la commutativité du diagramme

% foot, triangle
\diagram{\begin{tikzcd}[ampersand replacement=\&]
\& {}_{n'}P\otimes Q\arrow[dd,"{\handwritten{m\otimes \mathrm{id}_Q}}"]\\
\underline{\mathrm{Tor}}_1(P,Q)\arrow[ur]\arrow[dr] \&\\
\& {}_nP\otimes Q\quad\deleted{\handwritten{\unsure{m|+}}}\quad\deleted{\handwritten{\unsure{q|4}}}\,,
\end{tikzcd}}

% foot, last lines
où les flèches obliques sont les deux isomorphismes du type (1.3.4), relatifs à $n$ et $n'$ respectivement. Or cette compatibilité résulte de
\endgroup
% archive sheet number: 29

\scanpage{31}

% top, opening
Supposons $j$ et $j'$ ainsi choisis. Si $j,j'$ sont injectifs, on en déduit une biextension \inserted{canonique} \deleted{de} $(A,A')$, où $A=G/M$, $A'=G'/M'$. Dans le cas général, posons

% line struck above complexes
\deleted{\unsure{peut dire qu'on a une biextension de|on dira qu'on a une biextension de}}

% upper complexes
\[ A=\bigl[\,\underset{\text{deg. }-1}{M}\longrightarrow\underset{\text{deg }0}{G}\,\bigr],\qquad A'=[M'\mathbin{-}G'] \]

% triangles, caption and ends
d'où
\diagram{\begin{tikzcd}[ampersand replacement=\&,column sep=small]
\& A\arrow[dl,"1"'] \& \& \text{\unsure{complexe} de} \& \& A'\arrow[dl,"1"']\&\\
\mathbin{-}M\arrow[rr] \& \& G\mathbin{-}\arrow[ul] \& \text{chaînes ?} \& \mathbin{-}M'\arrow[rr,dash] \& \& G'\arrow[ul]
\end{tikzcd}}

% class, introduction
on va \unsure{définir|déduire} une \deleted{\unsure{biexten.|biext}} classe canonique

% class
\[\xi\in\mathrm{Ext}^{1}(A\overset{\mathrm L}{\otimes}A',G_m),\]

% middle, role
jouant le rôle d'une biextension de $(A,A')$ par $G_m$.

% left of gap
% slanted
\margin{\unsure{bien orientés|remonter|pour ainsi dire} ?}

% middle, unfinished phrase
Soit en effet \ldots

% lower, definitions
\deleted{Soit\unsure{ de}}\quad $\,{}_nA\overset{\text{déf}}{=}\deleted{\unsure{R}}\underline{\mathrm{Ext}}^{0}(\deleted{\unsure{M}}\mathbb{Z}/n\mathbb{Z},A)$, $\,{}_nA'=\underline{\mathrm{Ext}}^0(\mathbb{Z}/n\mathbb{Z},A')$

% lower, transition
on trouve, \unsure{grâce}, par \unsure{aux triangles|aux hypothèses} plus haut,

% lower, left margin
% slanted
\margin{donner aussi autre interprétat\unsure{°}. de $\,{}_nA$, via hom cobord si $nG=G$}

% lower, struck word and brace
\uline{\deleted{\unsure{Soit|si}}}
\[\left\{\begin{aligned}
 {}_nM&=0\quad\deleted{{}_nM'=0}\inserted{\unsure{\text{resp.}}}\\
 nG&=G\quad\text{p. ex.}\quad\deleted{\unsure{\text{par}|{}_nB}}\;nB=B,\quad nT=T\quad(\text{ce dernier résulte de }{}_nM=0
\end{aligned}\right.\]

% lower, torsion sequence
\[\unsure{\cdot|0}\longrightarrow\underbrace{{}_nG}_{D(M'_n),\,{}_nB),\,\text{si }nT=T}\longrightarrow{}_nA\longrightarrow M_n\longrightarrow 0\]

% lower, parallel sequence
et \unsure{on a de} $=$ type suite exacte
\[\cdot\mathbin{-}{}_nG'\mathbin{-}{}_nA'\longrightarrow M'_n\mathbin{-}0\]

% foot, passage text
sous les hyp. duales ; passant à la limite sur $n=\ell^\nu$, on a

% foot, two sequences
\[\left\{\begin{aligned}
\cdot&\longrightarrow T_\ell(G)\longrightarrow T_\ell(A)\longrightarrow M\otimes\mathbb Z_\ell\longrightarrow 0\\
\cdot&\mathbin{-}T_\ell(A')\mathbin{-}T_\ell(A')\mathbin{-}M'\otimes\mathbb Z_\ell\longrightarrow 0
\end{aligned}\right.\]
% archive sheet number: 30

\scanpage{32}
\begingroup\csname tagsleft@true\endcsname

% top left, two circled models
% slanted
\margin{\handwritten{\tikz[baseline=(c.base)]{\node[draw,ellipse,inner sep=1pt] (c) {$\phi^{(\ell)}$};}\ \tikz[baseline=(c.base)]{\node[draw,ellipse,inner sep=1pt] (c) {${}_{\ell^\nu}P$};}}}

% first formula
\[\deleted{\unsure{(}}\phi_{\handwritten{\xi}}^{(\ell)}\deleted{\unsure{)}}=(\phi_{\handwritten{\xi}}^{\ell^{\handwritten{\nu}}})_{\handwritten{\nu\geq}0}\,,\quad\text{où}\quad \phi_{\handwritten{\xi}}^{\ell^{\handwritten{\nu}}}:{}_{\ell^{\handwritten{\nu}}}P\;\mathsf{x}\;{}_{\ell^{\handwritten{\nu}}}Q\handwritten{\longrightarrow}G\,.\tag*{(2.2.4)}\]

% top, diagrams and systems
Les diagrammes\inserted{ commutatifs} déduits de (2.2.2) en y faisant $n=\ell^{\handwritten{\nu}}$, $n'=\ell^{\handwritten{\nu'}}$, pour $\handwritten{\nu'\geq\nu}$, montrent que $\phi_{\handwritten{\xi}}^{(\ell)}$ peut être interr̂été comme un accouplement de \deleted{pro-groupes} systèmes projectifs de groupes

% second formula
\[\phi_{\handwritten{\xi}}^{(\ell)}:\ T_\ell(P)\mathsf{x}\,T_\ell(Q)\handwritten{\longrightarrow}T_\ell(G)\,,\tag*{(2.2.5)}\]

% upper middle, introduction
où, pour un \deleted{\unsure{G}}groupe $M$ de $T$ et un nombre premier $\ell$, on désigne par

% projective system
\[T_\ell(M)=({}_{\ell^{\handwritten{\nu}}}M)_{\handwritten{\nu\geq}0}\]

% upper middle, name
le système p\deleted{\unsure{(}}rojectif des $\,{}_{\ell^{\handwritten{\nu}}}M$. L'accouplement \deleted{\unsure{$\phi$|à}} (2.2.5) est encore appelé \uline{l'accouplement} $\ell$-\uline{adique associé à} $\handwritten{\xi}$ (ou à la biextension $E$ de classe $\handwritten{\xi}$). Pour $\handwritten{\xi}$ variable, on obtient un homomorphisme canonique

% third formula
\[\mathrm{Biext}^{\handwritten{1}}(P,Q;G)\handwritten{\longrightarrow}\mathrm{Hom}(T_\ell(P)\otimes T_\ell(Q),T_\ell(G))\,.\tag*{(2.2.6)}\]

% middle, functoriality
On notera que les homomorphismes (2.1.2) sont \uline{fonctoriels} en $P,Q,G$ \uuline{(pour $n$} fixé), d'où on déduit qu'il en est de même de (2.2.6).

% section 2.3
2.3.\qquad \deleted{On peut regarder} La compatibilité (2.2.2) permet de retrouver $\phi_{\handwritten{\xi}}^n$ connaissant $\phi_{\handwritten{\xi}}^{n'}$, pourvu que $\,{}_{n'}Q\handwritten{\xrightarrow{\mathrm m}}{}_nQ$ \handwritten{soit} un épimorphisme, -- \deleted{(NB \unsure{il}} condition satisfaite, comme il est bien connu, si $nQ=Q$. En effet, comme $\,{}_{n'}P\handwritten{\xrightarrow{\mathrm m}}{}_nP$ est aussi un épimorphisme, $\phi_{\handwritten{\xi}}^n$ se déduit \deleted{\unsure{de}}alors de $\phi_{\handwritten{\xi}}^{n\handwritten{'}}$ par passage au quotient. On peut cependant \inserted{\handwritten{donner}} une compatibilité légèrement plus précise que (2.2.2), qui permet d\deleted{e}'exprimer $\phi^n$ en termes de $\phi^{n'}$ sans aucune restriction sur $Q$\handwritten{,} par la commutativité de

% lower square
\diagram{\begin{tikzcd}[ampersand replacement=\&,column sep=small]
\&{}_{n'}P\arrow[d,"\handwritten{m}"]\&\mathsf x\&{}_nQ\arrow[d,"\handwritten{\mathrm{id}}"]\arrow[rr,"{\phi_{\handwritten{\xi}}^{n'}\handwritten{|}\,{}_{\handwritten{n'}}P\;\mathsf{x}\;{}_nQ}"]\&\&{}_nG\arrow[d,"\handwritten{\mathrm{id}}"]\\
\text{(2.3.1)}\&{}_{\deleted{\unsure{m}}n}P\&\mathsf x\&{}_nQ\arrow[rr,"{\phi_{\handwritten{\xi}}^n}"]\&\&{}_nG\,,
\end{tikzcd}}

% under square
en formules

% penultimate formula
\[\phi_{\handwritten{\xi}}^n(mx,y)=\phi_{\handwritten{\xi}}^{n'}(x,y)\qquad(x\in{}_{n'}P(S),\ y\in{}_nQ(S)).\tag*{(2.3.2)}\]

% last lines
Cette formule implique évidemment la commutativité de (2.2.2), qui s'écrit en effet

% last formula
\[\phi_{\handwritten{\xi}}^n(mx,my)=m\,\phi_{\handwritten{\xi}}^{n'}(x,y)\qquad(x\handwritten{\in}{}_{n'}P(S),\ y\handwritten{\in}{}_{n'}P(S)),\tag*{\handwritten{(2.3.3)}}\]
\endgroup
% archive sheet number: 31

\scanpage{33}

% top, cancelled opening
\deleted{\unsure{Les inclusions|Les injections}}

% top, maps
on a \uline{loc.} des \unsure{hom|$\mathbb G_m$}
\[{}_nA\longrightarrow A,\quad {}_nA'\mathbin{-}A'\]
d'où \uline{loc} des \unsure{maps|composés|ces morphismes}

% top, Ext triangle
\diagram{\begin{tikzcd}[ampersand replacement=\&]
\mathrm{Ext}^{1}(A\overset{\mathrm L}{\otimes}A,G_m)\arrow[r]\arrow[dr,dashed] \& \mathrm{Ext}^{1}({}_nA\overset{\mathrm L}{\otimes}{}_nA',G_m)\arrow[d]\\
\& \mathrm{Hom}({}_nA\otimes{}_nA',G_m)
\end{tikzcd}}

% upper middle, gluing
qui se recollent, donc sont définis \unsure{glbᵗ|glblᵗ|globalt}. Donc l'élément $\xi$ envisagé plus haut nous donne

% upper middle, pairing
\[\varphi_n^\xi:{}_nA\otimes{}_nA'\longrightarrow G_m\]

% middle, limit
et par passage à la limite
\[T_\ell(A)\times T_\ell(A')\longrightarrow T_\ell(G_m)\overset{\text{déf}}{=}\mathbb Z_\ell(1).\]

% proposition, opening
\uline{Prop} Dans cet accouplement, $D(\unsure{M'_n|{}_nM_n|M_n})$ est orthogonal à $\,{}_nG'$, \deleted{\unsure{et}} (a un sens si $\,{}_nM'=0$) et \deleted{$D(M_n$} $\,{}_nG$ est orth. à $D(M_n)$ (a un sens si $\,{}_nM=0$). \deleted{Les} accouplements induits \deleted{$\unsure{{}_nA/}$}

% proposition, first pairing and lower identifications
\diagram{\begin{tikzcd}[ampersand replacement=\&,column sep=small]
{}_nA'/{}_nG'\&\times\&D(M'_n)\&{}\longrightarrow G_m\\
M'_n\arrow[u,dash]\&\&{}_nT\unsure{{}|'}\arrow[u,equal]\arrow[ull,bend right=56,dash]\&
\end{tikzcd}}

% proposition, next pairing
et $\underset{{}_nT}{\underset{\Vert}{D(M_n)}}\times\underset{M_n}{\underset{|}{{}_nA/{}_nG}}\mathbin{-}G_m$, ainsi que

% proposition, last factors
\deleted{\unsure{et}}, $\underset{{}_nB}{\underset{\Vert}{{}_nG/{}_nT}}\deleted{\unsure{{}_nT}}\times\underset{{}_nB'}{\underset{|}{{}_nG'/{}_nT'}}$, sont ceux qu'on devine.

% foot, corollary
\uline{Corollaire} Si l'accouplement entre $\,{}_nB$ et $\,{}_nB'$ est une dualité parfaite (cas où $T_n=T'_n=0$) \deleted{\unsure{alors}} et si $\,{}_nM=0$, $\,{}_nM'=0$, $nB=B$, $nB'=B'$, alors l'accouplement de $\,{}_nA$ et $\,{}_nA'$ est une dualité parfaite.
% archive sheet number: 32

\scanpage{34}
\begingroup\csname tagsleft@true\endcsname

% heading
2.\quad \uline{Accouplements définis par une biextension}.

% section 2.1
2.1.\quad Soit $n$ un entier $\handwritten{>}0$, et supposons qu'on a

% first assumption
\[nP=P\,.\tag*{(2.1\deleted{)}.1)}\]

% composite, lead
Considérons alors l'homomorphisme composé

% composite
\diagram{\begin{tikzcd}[ampersand replacement=\&,column sep=small]
\text{(2.1.2)}\&\mathrm{Biext}^{\handwritten{1}}(P,Q;G)\arrow[r]\arrow[rr,bend left=18,"{\handwritten{\varphi^n}}"]\&\mathrm{Biext}^{\handwritten{1}}(P,{}_nQ;G)\arrow[r,"\handwritten{\sim}"]\&\mathrm{Hom}({}_nP\handwritten{\otimes}{}_nQ;G)\,,
\end{tikzcd}}

% restriction and number
où la première flèche est la restriction à $(P,{}_nQ)$, et la deuxième est l'isomorphisme (1.3.\handwritten{6}). Il permet, à tout élément
\[\inserted{\handwritten{\xi\in}\mathrm{Biext}(P,Q;G)}\,,\tag*{(2.1.3)}\]
a fortiori à toute biextension\inserted{ $E$} de $(P,Q)$ par $G$, de faire correspondre un\deleted{e forme homomorphis\unsure{m|me}} accouplement bilinéaire

% margin beside pairing
\margin{\handwritten{\tikz[baseline=(c.base)]{\node[draw,ellipse,inner sep=1pt] (c) {$\varphi$};}}}

% finite pairing
\[\phi_{\handwritten{\xi}}\ \text{ou}\ \phi_{\handwritten{\xi}}^n:\ {}_nP\;\mathsf{x}\;{}_nQ\handwritten{\longrightarrow}G\,,\tag*{(2.1.4)}\]

% name of pairing
qu'on appelle l'accouplement \uline{associé à} $\handwritten{\xi}$ (ou à la biextension $E$).

% section 2.2
2.2.\quad Soit \deleted{\unsure{$n$}} $n'$ un multiple de $n$,
\[n'=mn\,,\]
et supposons qu'on ait la relatiòn (plus forte que (2.1.1))

% second assumption
\[n'P=P\quad\text{i.e.}\quad mnP=P\,,\tag*{(2.2.1)}\]

% new pairing, introduction
d'où comme en (2.1.4), pour toute $\handwritten{\xi}\in\mathrm{Biext}(P,Q;G)$, un\deleted{e} accouplement

% new pairing
\[\phi_{\handwritten{\xi}}^{n'}:\ {}_{n'}P\;\mathsf{x}\;{}_{n'}Q\handwritten{\longrightarrow}G\,.\]

% revised introduction to square
\deleted{Les relations entre cet accouplement et}\inserted{\handwritten{Il est lié à}} l'accouplement $\phi_{\handwritten{\xi}}^n$ \deleted{est donné} par la commutativité du diagramme suivant (\deleted{qu'on laisse au lecteur le soin d'établir) :}\inserted{\handwritten{cf. 2.3 ci-dessous pour une démonstration) ;}}

% lower square
\diagram{\begin{tikzcd}[ampersand replacement=\&,column sep=small]
\&{}_{n'}P\arrow[d,"\handwritten{m}"']\&\handwritten{\times}\&\deleted{\unsure{{}_nQ}}\,\handwritten{{}_{n'}Q}\deleted{\unsure{Q|n}}\arrow[d,"\handwritten{m}"]\arrow[rr,"{\phi_{\handwritten{\xi}}^{n'}}"]\&\&{}_{\handwritten{n'}}G\arrow[d,"\handwritten{m}"]\\
\text{(2.2.2)}\&\deleted{\unsure{P}}\,\handwritten{{}_nP}\&\deleted{\unsure{\times}}\handwritten{\times}\&{}_nQ\arrow[rr,"{\phi_{\handwritten{\xi}}^n}"]\&\&{}_{\handwritten{n}}G\,\handwritten{,}
\end{tikzcd}}

% lower inserted explanation
\inserted{\handwritten{où les flèches verticales sont induites par la multiplication par $m=n'/n$.}}

% lower paragraph
Notons maintenant que moyennant (2.1.1), on peut prendre pour $n'$ n'importe quelle puissance de $n$. En particulier, soit $\ell$ un nombre premier, et supposons qu\deleted{\unsure{'}}e\deleted{\unsure{,}} l'on ait

% lower left ell
\margin{\handwritten{\tikz[baseline=(c.base)]{\node[draw,ellipse,inner sep=1pt] (c) {$\ell$};}}}

% third assumption
\[\ell P=P\,,\tag*{(2.2.3)}\]

% last line
et considérons le système des accouplements
\endgroup
% archive sheet number: 33

\scanpage{35}

% upper right, underlined heading
\uline{Picard et dualité}
% archive sheet number: 34

\scanpage{36}

% left of first two bands
% margin note written sideways
\margin{$\unsure{\rightleftarrows|{\ast\atop\shortmid}|\text{et}}$ Cohomologie discrète}

% three-column cohomology table, read across horizontal bands
\diagram{\resizebox{\linewidth}{!}{$\begin{array}{@{}r@{}c|c|c}
& H^0(X,\unsure{\zeta|\xi|\mathfrak z}_\ell)=U_\ell &
\begin{gathered}H^1(X,\zeta_\ell)=P^{(\ell)}\quad\text{éléments de la jacobienne}\\[-1pt]\text{$P$ de $X$ annulés par $\ell$}\end{gathered}
& H^2(X,\zeta_\ell)\simeq\mathbb Z/\ell\\[5pt]
&\begin{aligned}
H^2(X,\mathbb Z/\ell)&=H^2(X,\zeta_\ell)\otimes_{\mathbb Z/\ell}\mathrm{Hom}(U_\ell,\mathbb Z/\ell)\\
 &=\mathbb Z_\ell\otimes_{\mathbb Z/\ell}\mathrm{Hom}(U_\ell,\mathbb Z/\ell)\\
 &=\deleted{\mathbb Z/\ell}\,\mathrm{Hom}(U_\ell,\mathbb Z/\ell)
\end{aligned}
&
\begin{aligned}
H^1(X,\mathbb Z/\ell)&=H^1(X,\zeta_\ell)\otimes_{\mathbb Z/\ell}\mathrm{Hom}(\deleted{\unsure{H}}U_\ell,\mathbb Z/\ell)\\
&\deleted{=P^{(\ell)}\otimes_{\text{''}}\mathrm{Hom}(\zeta_\ell,H)}\\
&=P^{(\ell)}\otimes_{\mathbb Z/\ell}\mathrm{Hom}(\deleted{\unsure{H}}U_\ell,\mathbb Z/\ell)
\end{aligned}
& H^0(X,\mathbb Z/\ell)\simeq\mathbb Z/\ell\\
\smash{\left\{\rule[-10ex]{0pt}{25ex}\right.}&&\begin{gathered}\text{Or}\quad P^{(\ell)}\times P^{(\ell)}\longrightarrow U_\ell\quad\text{\unsure{en vertu}}\\
\boxed{\text{de}\ \boxed{\begin{gathered}\text{l'autodualité de $P$ ---}\\\text{\uline{et} la dualité de \unsure{Kummer-Serre|Kummer Serre} --}\end{gathered}}}\end{gathered}&\\ \hline
& H^0(X,\zeta_p)=0&
\begin{gathered} H^{\deleted{\unsure{2}}1}(X,\zeta_p)=P^{(p)}=\text{éléments d'ordre $p$}\\\text{de la jacobienne}\end{gathered}&H^2(X,\zeta_p)\simeq\mathbb Z/p\\
& H^2(X,\mathbb Z/p)\simeq H^2(X,\mathbb Z/p)=0&
\begin{gathered} H^1(X,\mathbb Z/p)\simeq\text{éléments de $H^1(X,\underline{O}_X)$ fixes}\\\text{par $F$}\\\boxed{\text{dualité de Serre - Cartier}}\end{gathered}
&H^0(X,\mathbb Z/p)\simeq\mathbb Z/p\\ \hline
& H^0(X,i_p)=0&
H^1(X,\deleted{\unsure{O|\mathfrak z}}i_p)\simeq {}_F H^1(X,\underline{O}_X)^{\deleted{\unsure{F}}}
&H^2(X,i_p)=H^1(X,\underline{O}_X)_F
\end{array}$}}
% archive sheet number: 35

\scanpage{37}

% top, cohomology
\[H^1(X,\underbrace{\underline{\mathrm{Ext}}^1(f^*(B),G_{m\,X})})\]

% upper triangle and upper identification
\diagram{\begin{tikzcd}[ampersand replacement=\&]
\unsure{?}\quad\mathrm{Ext}_X(\overbrace{f_{\unsure{!|*|{}}}^*(B)}^{A},G_{m\,X})\arrow[r,dash,phantom,"\simeq"] \& \mathrm{Ext}_S(B,f_!(G_{m\,X}))\\
H^p(X,\underline{\mathrm{Ext}}_X^q(A,G_{m\,X}))\arrow[u,Rightarrow]\&
\end{tikzcd}}

% first horizontal rule
\hrule

% middle, brace over list
\[\left\{\begin{array}{l}
\text{\unsure{Suite exacte|Suites exactes} de Cohomologie. Critère de \unsure{trivialité|representabilité} ;}\\
\text{\unsure{Frobeniuseries}}\\
\text{Dualité de Cartier \quad [pour Pic ---]}\\
\text{Th. qualitative des V.A. \quad (Th. de rigidité.)}\\
\text{Propriétés \unsure{générales|formelles} des groupes ---}\\
\text{Propriétés différentielles \deleted{des} groupes et fibrés}\\
\text{\unsure{Cas particuliers|Cas particulier} : \uuline{fibrés linéaires}, symplectiques, orthogonaux}\\
\hfill\text{etc ---}
\end{array}\right.\]

% middle, detached question
Sans doute pour les représentations \unsure{lin.} de \unsure{grp} ??

% second horizontal rule
\hrule

% bottom, cancelled Hom
\[\deleted{\underline{\mathrm{Hom}}(G_a,G_m)}\]

% bottom, last symbols
\[\mathbb Z\unsure{/n|\text{ sur }|\longmapsto|\longrightarrow}\mathbb Z\]
% archive sheet number: 36
